\section*{अध्याय \ref{ch:b2:speciation} --- जाति-उद्भव और महाविकास}

\begin{solution}{exo:b2:speciation:1}
कोई जाति आपस में प्रजनन करती प्राकृतिक समष्टियों का समूह है जो
दूसरे ऐसे समूहों से जननिक रूप से विलगित हो। वह
अलैंगिक वंशक्रमों के लिए चूकती है (जीवाणु, ब्डेलॉइड रोटिफ़र, सिंहपर्णी:
यानी आकारिकीय या जातिवृत्तीय संकल्पना, एक जैसे रूपों के गुच्छ या
पहचान-योग्य वंशक्रम); जीवाश्मों के लिए (आकारिकीय संकल्पना, क्योंकि
कोई संकरण किया ही नहीं जा सकता); और ऐसी भिन्नक्षेत्रीय समष्टियों के लिए जो कभी
मिलती ही नहीं (पारिस्थितिक या जातिवृत्तीय संकल्पना, क्योंकि आपसी प्रजनन
जाँचा नहीं जा सकता) --- और वह संकरित होते बलूतों से तथा
वलय-जातियों से धुँधली होती है।
\end{solution}

\begin{solution}{exo:b2:speciation:2}
भिन्न महीनों में मेंढक: पूर्वयुग्मनजी (कालिक)। खच्चर: पश्चयुग्मनजी
(संकर-बंध्यता)। ऐसा शुक्राणु जो बँध न सके: पूर्वयुग्मनजी (युग्मकीय)।
बंध्य \omterm{def:b2:angiosperm-reproduction:seed}{बीज} वाले ओजस्वी संकर: पश्चयुग्मनजी (संकर-बंध्यता,
या संकर-भंग यदि वह विफलता अगली पीढ़ी में हो)।
भिन्न गान वाले झींगुर: पूर्वयुग्मनजी (व्यवहारिक)।
\end{solution}

\begin{solution}{exo:b2:speciation:3}
भिन्नक्षेत्रीय: कोई बाधा किसी परास को बाँटती है (पनामा स्थलडमरूमध्य के दोनों
पार्श्वों पर चटकाती झींगा-जातियाँ)। परिधीय: उस परास के परे कुछ संस्थापक
(किसी मुख्यभूमि पूर्वज से डार्विन के फ़िंच)। समक्षेत्रीय:
कोई बाधा नहीं (नागफनी तथा सेब पर वह सेब-कीट मक्खी; कोई
\omterm{def:b2:genome-diversification:polyploidy}{विषमबहुगुणित} गेहूँ)। परिक्षेत्रीय: किसी प्रवणता के साथ, बीच में किसी
\omterm{prop:b2:speciation:reinforcement}{संकर-क्षेत्र} समेत (खान-मलबे पर तथा उससे बाहर की घासें, जो
भिन्न समयों पर फूलती हैं)।
\end{solution}

\begin{solution}{exo:b2:speciation:4}
जाति 1 सदा जीतती है; जाति 2 सदा जीतती है; स्थिर सहअस्तित्व;
और आरंभिक संख्याओं के अनुसार कोई भी जीतती है। सहअस्तित्व तब स्थिर है
जब हर जाति दूसरी पर उसकी धारण-क्षमता पर आक्रमण कर सके,
यानी जब हर जाति स्वयं को उससे अधिक सीमित करे जितना वह
दूसरी को करे ($K_1 > \alpha K_2$ तथा $K_2 > \beta K_1$)।
\end{solution}

\begin{solution}{exo:b2:speciation:5}
$\alpha K_2 = 400 < K_1 = 1000$: जाति 1 आक्रमण करती है; $\beta K_1 = 600
< K_2 = 800$: जाति 2 आक्रमण करती है। $N_1^{*} = (1000 - 400)/(1 - 0.3) =
857$, $N_2^{*} = (800 - 600)/0.7 = 286$; और उस साम्य पर जाति
1 अपने $K$ से $\alpha N_2^{*} = 143$ नीचे है और जाति 2 $\beta
N_1^{*} = 514$।
\end{solution}

\begin{solution}{exo:b2:speciation:6}
$\alpha K_2 = 1200 > K_1$: जाति 1 जाति 2 पर उसकी
धारण-क्षमता पर आक्रमण नहीं कर सकती; $\beta K_1 = 300 < K_2$: जाति 2
जाति 1 पर आक्रमण कर सकती है। जाति 2 सदा जीतती है, क्योंकि उसके व्यष्टि
जाति 1 को उससे अधिक चोट पहुँचाते हैं जितना जाति 1 के अपने, जबकि वह स्वयं मुश्किल से
प्रभावित होती है। $\alpha = \beta = 1.5$ के साथ: $\alpha K_2 = 1200 > 1000$
और $\beta K_1 = 1500 > 800$: कोई भी दूसरी पर आक्रमण नहीं कर सकती, हर एक
जम जाने पर दूसरी को अपवर्जित कर देती है, और संस्थापक जीतता है।
\end{solution}

\begin{solution}{exo:b2:speciation:7}
चतुर्गुणित $4n$ $2n$ युग्मक बनाता है; \omterm{def:b2:meiosis-heredity:meiosis}{द्विगुणित} के $n$
युग्मकों से संकरित होकर वह त्रिगुणित $3n$ देता है, जिनका अर्धसूत्रण तीन
समुच्चय युग्मित नहीं कर सकता और अगुणितेतर, अधिकतर अजीवनक्षम युग्मक देता है। वह
चतुर्गुणित केवल स्वयं से प्रजनन कर सकता है (या स्व-परागण से): यानी
एक ही चरण में विलगित। वह आम तौर पर इसलिए चूक जाता है कि वह अकेला है --- उसके एकमात्र
संभावित साथी \omterm{def:b2:meiosis-heredity:meiosis}{द्विगुणित} हैं और हर ऐसा संकरण व्यर्थ जाता है ---
और वह दबा दिया जाता है; वह तब टिकता है जब वह स्व-परागण करे, कायिक रूप से प्रवर्धित हो,
या बार-बार उपजे।
\end{solution}

\begin{solution}{exo:b2:speciation:8}
$0.01\times 10\times 10 = 1$ असंगति; और हर एक 30 के साथ, $0.01\times
900 = 9$। विलगन विचलन-काल के वर्ग की तरह बढ़ता है: तीन गुना
लंबे समय तक विचलित दो समष्टियाँ नौ गुना
विलगित हैं --- यानी वह हिमगोला।
\end{solution}

\begin{solution}{exo:b2:speciation:9}
$w = 2\sqrt{8/0.05} = \qty{25}{km}$; $w = 0.3\sqrt{8/0.4} =
\qty{1.3}{km}$। प्रबलन दूसरे में अधिक संभव है: संकर
ज़ोर से अनुपयुक्त हैं, इसलिए जो मादा दूसरी क़िस्म से इनकार करे उसे
बहुत मिलता है, और वह क्षेत्र इतना सँकरा है कि वे दोनों रूप उस इनकार के वरे जाने भर
अक्सर मिलें।
\end{solution}

\begin{solution}{exo:b2:speciation:10}
$\alpha K_2 = 450 < 500$, इसलिए हर एक दूसरी पर आक्रमण करती है: वे सह-अस्तित्व में रहती हैं,
हर एक $N^{*} = (500 - 450)/(1 - 0.81) = 263$ पर --- मुश्किल से, अपवर्जन के
किनारे के पास, और $K$ में कोई छोटा परिवर्तन उसे लुढ़का देता।
$\alpha = \beta = 0.4$ के साथ: हर एक $N^{*} = (500 - 200)/(1 - 0.16) = 357$।
विस्थापन ने हर समरेख को दूसरे के समरेख से दूर घुमा दिया ($K/\alpha$
बाहर की ओर खिसका), इसलिए वह कटान अक्षों से दूर हट गया और हर
जाति अपने ही $K$ के अधिक पास बैठती है: यानी अधिक सुरक्षित सहअस्तित्व,
और हर एक की अधिक संख्या।
\end{solution}

\begin{solution}{exo:b2:speciation:11}
$R = h^{2}S = 0.7\times(-1) = \qty{-0.7}{mm}$: यानी जैसा देखा गया।
1977 के सूखे में, जब कोई बड़ा भू-फ़िंच मौजूद नहीं था, वही मध्यम
फ़िंच \emph{बड़ी} चोंचों के लिए वरे गए थे, क्योंकि बड़े \omterm{def:b2:angiosperm-reproduction:seed}{बीज} ही
बचा भोजन थे; 2004 में बड़े \omterm{def:b2:angiosperm-reproduction:seed}{बीज} उस
हमलावर ने ले लिए, और जो मध्यम फ़िंच उनके लिए होड़ करते थे वे हारे: यानी
वरण की दिशा उस होड़ी ने बाँधी।
\end{solution}

\begin{solution}{exo:b2:speciation:12}
संयोग से बनीं: वह दो-स्थल प्रतिरूप दिखाता है कि विलगन दूसरे कारणों से
नियत हुए प्रतिस्थापनों के किसी उपोत्पाद के रूप में, या अपवाह से, उपजता है,
और भूगोल --- कोई बाधा, कोई संस्थापन --- वह अलगाव देता है
जो उस संयोग को जमा होने देता है। वरण से रखी गईं: जहाँ वे
रूप मिलें, वहाँ अनुपयुक्त संकरों के विरुद्ध वरण पूर्वयुग्मनजी
बाधाएँ बनाता है (प्रबलन), और होड़ लक्षण इस तरह विस्थापित करती है
कि वे दोनों सह-अस्तित्व में रह सकें। वह उक्ति विदारक वरण से \omterm{prop:b2:speciation:modes}{समक्षेत्रीय जाति-उद्भव}
छोड़ देती है, जहाँ वरण उन जातियों को बनाता भी है और रखता भी,
और \omterm{def:b2:genome-diversification:polyploidy}{बहुगुणिता} भी, जहाँ वह संयोग ही पूरी कहानी है।
\end{solution}

\begin{solution}{pb:b2:speciation:1}
\textbf{1.} जाति 1: $N_1 + 0.9N_2 = 1200$, अंतःखंड $N_1 = 1200$
तथा $N_2 = 1333$। जाति 2: $N_2 + 0.3N_1 = 400$, अंतःखंड $N_2 =
400$ तथा $N_1 = 1333$।
\textbf{2.} $\beta K_1 = 360 < K_2 = 400$: जाति 2 आक्रमण करती है;
$\alpha K_2 = 360 < K_1 = 1200$: जाति 1 आक्रमण करती है। यानी स्थिर
सहअस्तित्व।
\textbf{3.} $N_1^{*} = (1200 - 360)/(1 - 0.27) = 1151$; $N_2^{*} =
(400 - 360)/0.73 = 55$।
\textbf{4.} $\mathrm{d}N_1/\mathrm{d}t = 0.5\times 1200\,(1 - 1218/1200)
= -9$ प्रति वर्ष: यानी वह स्थानीय, अपनी धारण-क्षमता पर, उन नवागंतुकों से
ह्रास में धकेल दिया जाता है।
\textbf{5.} $K_1 = 600$, $K_2 = 200$: $N_1^{*} = (600 - 180)/0.73 =
575$, $N_2^{*} = (200 - 180)/0.73 = 27$: दोनों आधे हो जाते हैं, और वह हमलावर,
27 पक्षियों पर, विलोपन के पास है --- यानी सूखा ही वह समय है जब
होड़ काटती है।
\textbf{6.} वह हमलावर बड़ी चोंच वाला बड़ा पक्षी है: वह
वे बड़े \omterm{def:b2:angiosperm-reproduction:seed}{बीज} लेता है जिन्हें वह स्थानीय भी काम में लेता है, और उन्हें तेज़ लेता है,
जबकि उस स्थानीय के छोटे \omterm{def:b2:angiosperm-reproduction:seed}{बीज} उसके थोड़े ही काम के हैं।
\textbf{7.} $S = 9.0 - 9.5 = \qty{-0.5}{mm}$; $R = 0.7\times(-0.5) =
\qty{-0.35}{mm}$।
\textbf{8.} $9.5 - 0.35 = \qty{9.15}{mm}$।
\textbf{9.} $N_1^{*} = (1200 - 200)/(1 - 0.15) = 1176$; $N_2^{*} =
(400 - 360)/0.85 = 47$।
\textbf{10.} जाति 1 के समरेख का $N_2$ अंतःखंड
$K_1/0.9 = 1333$ से चढ़कर $K_1/0.5 = 2400$ हो गया: वह रेखा उस
हमलावर की रेखा से दूर घूम गई, वह स्थानीय अब हर हमलावर से कम प्रभावित है,
और वह साम्य उस स्थानीय के $K$ की ओर खिसक गया।
\textbf{11.} $2/0.35 \approx 6$ वर्ष। वह इसलिए नहीं होता कि इस तरह का वरण
केवल सूखे वर्षों में होता है, गीले वर्षों में उलट जाता है
जब छोटे \omterm{def:b2:angiosperm-reproduction:seed}{बीज} बहुत हों (या जब वह हमलावर दुर्लभ हो), और
क्योंकि वह योगात्मक प्रसरण चुक जाता।
\textbf{12.} बीस लाख वर्ष लगभग $10^{6}$ पीढ़ियाँ हैं; और प्रति पीढ़ी
\qty{0.35}{mm} की कोई खिसक उनके हज़ारवें भाग भर भी बनी रहे तो
वह चोंच सैकड़ों मिलीमीटर बदल जाती। वरण
मज़बूत पर अस्थिर है; बनी रहती दिशात्मक बदल दुर्लभ
अपवाद है, और भूवैज्ञानिक समय पर शुद्ध परिवर्तन बड़े विपरीत प्रसंगों का
कोई नन्हा शेष है।
\textbf{13.} गान, जो पिता से सीखा जाता है और मादाएँ जिसे
साथी-चुनाव में काम में लेती हैं, और स्वयं चोंच का आकार भी संकेत के रूप में। मध्यवर्ती चोंचों वाले
संकर न बड़े \omterm{def:b2:angiosperm-reproduction:seed}{बीज} सबसे अच्छे सँभालते हैं न छोटे,
और सूखों में सबसे पहले भूखे मरते हैं।
\textbf{14.} $w = 0.5\sqrt{8/0.2} = \qty{3.2}{km}$।
\textbf{15.} वह द्वीप उस क्षेत्र से, जितना वह होता, सँकरा है: कोई
स्थानिक क्लाइन बन ही नहीं सकता; और संकरण, जहाँ हो, वहाँ
पूरी समष्टि को प्रभावित करता है, तथा उन दोनों को केवल समजातीय प्रजनन ही
अलग रखता है।
\textbf{16.} जो मादाएँ केवल अपने ही गान तथा चोंच वाले नरों से प्रजनन करें वे
संकरण करने वालियों से अधिक नाती-पोते छोड़ती हैं; और वह
अभिरुचि तथा वह संकेत (गान, चोंच) वहाँ अधिक भिन्न हो जाते हैं जहाँ वे
दोनों जातियाँ साथ रहती हैं, बनिस्बत जहाँ वे अकेली रहती हैं।
\textbf{17.} $0.005\times 15\times 15 = 1.1$ असंगतियाँ: यानी
औसतन एक, और संकर शायद घटे हुए पर बंध्य नहीं। जैविक कसौटी से अभी
जातियाँ नहीं; पर रास्ते पर।
\textbf{18.} जैविक: अभी नहीं (संकर लगभग अनुकूल)।
आकारिकीय: शायद, यदि उनकी चोंचें या पंखावरण दिखने लायक
विचलित हुए हों। जातिवृत्तीय: हाँ, यदि हर समष्टि का कोई नियत
पहचान-भेद हो। वे संकल्पनाएँ ठीक जाति-उद्भव की प्रक्रिया भर असहमत होती हैं,
और वही उन्हें उपयोगी बनाता है।
\textbf{19.} प्रति जाति प्रति वर्ष $1/(5\times 10^{6}) = 2\times 10^{-7}$;
और $10^{7}$ जातियों के लिए, प्रति वर्ष दो विलोपन।
\textbf{20.} प्रति वर्ष 200 से 2000 जातियाँ।
\textbf{21.} $0.75\times 10^{7}/10^{5} = 75$ जातियाँ प्रति वर्ष --- यानी
निचले सिरे पर वर्तमान दर से नीची। वर्तमान प्रसंग, यदि
बना रहे, तो कोई \omterm{prop:b2:speciation:tempo}{सामूहिक विलोपन} है जो उन पाँच बड़ों से तेज़
चल रहा है।
\textbf{22.} हारने वालों से ख़ाली हुए निकेत होड़ियों से
ख़ाली हैं, इसलिए कोई भी बचा हुआ रूपांतर जो किसी एक को काम में ले सके अनुकूल है,
और वे बचे हुए उनमें विविध हो जाते हैं।
\textbf{23.} \omterm{prop:b2:speciation:tempo}{विरामित साम्य}; और किसी छोटी विलगित समष्टि में
भिन्नक्षेत्रीय (परिधीय) जाति-उद्भव, जिसका वंशज फिर
मुख्य समष्टि के परास में फैल जाता है।
\textbf{24.} $1.04^{n} = 2$: $n = \ln 2/\ln 1.04 = 18$ पीढ़ियाँ।
वह अभिलेख करोड़ों वर्ष इसलिए दिखाता है कि वरण उलट जाता है,
\omterm{prop:b2:population-genetics:kinds}{स्थायीकारी वरण} अधिकांश समय हावी रहता है, और शुद्ध परिवर्तन
बड़े विपरीत प्रसंगों का छोटा अंतर है।
\textbf{25.} विस्थापन से पहले: $N_1^{*} =
1151$, $N_2^{*} = 55$ पर स्थिर सहअस्तित्व; बाद: 1176 तथा 47। एक पीढ़ी में चोंच की खिसक
\qty{-0.35}{mm}। \omterm{prop:b2:speciation:reinforcement}{संकर-क्षेत्र} चौड़ाई \qty{3.2}{km}।
\end{solution}
